%% =====================================================================
%% APPENDIX: Chapter 3 — Supplementary Methodology Materials
%% Content removed from Chapter 3 during examiner feedback and audit.
%% Reason: Focus Ch.3 on core methodology; move implementation details,
%% tool comparisons, and system architecture here for reference.
%% =====================================================================
\setchaptercolor{3}

\label{app:ch3_supplementary}

This appendix contains supplementary materials that support the research methodology presented in Chapter~\ref{chap:methodology}. These materials were separated from the main chapter to maintain focus on core methodological decisions while preserving detailed implementation specifications, tool comparisons, and system architecture documentation for reference.

%% =====================================================================
\section{Extended Epistemological Justification}
\label{sec:appx_epistemology}
%% =====================================================================

This research adopts a \textbf{pragmatist paradigm}~\cite{creswell2018research} as its philosophical foundation, situated within a predominantly \textbf{quantitative} approach augmented by design science elements. Pragmatism holds that knowledge claims arise from actions, situations, and consequences rather than from antecedent conditions alone, making it particularly well-suited to applied research that must bridge computational experimentation with clinical decision-making~\cite{creswell2018research}. The ontological position underpinning this study is one of \textbf{critical realism}: biomedical signals such as electroencephalographic recordings possess objective, measurable properties (voltage amplitudes, spectral power densities, temporal dynamics), yet the clinical interpretation of those signals is inherently context-dependent, shaped by diagnostic conventions, institutional protocols, and the experience of individual clinicians. This dual nature---objective measurement combined with socially situated interpretation---motivates the epistemological stance that knowledge in biomedical AI emerges from the convergence of quantitative empirical measurement and qualitative expert assessment, rather than from either source in isolation. Could a simpler design have sufficed? A purely positivist study could report accuracy; a purely interpretivist study could report clinician perceptions. Neither could do both with the traceability this dissertation demands.

The paradigm selection rationale is as follows:
\begin{itemize}
    \item \textbf{Why pragmatism:} Accommodates quantitative metrics (accuracy, F1, AUC), qualitative expert assessment, and constructive artifact evaluation within a single framework
    \item \textbf{Why not post-positivism:} Cannot recognize the framework artifact itself as a scholarly contribution
    \item \textbf{Why not interpretivism:} Lacks mechanisms for controlled DL model comparison
\end{itemize}


%% =====================================================================
\section{Research Timeline}
\label{sec:appx_timeline}
%% =====================================================================

Figure~\ref{fig:research_timeline_appx} presents the 26-week research timeline with four milestones, showing overlapping phases that indicate parallel work streams.
\needspace{20\baselineskip}
\begin{figure}[!htbp]
    \centering
    \resizebox{0.97\textwidth}{!}{%
    \begin{tikzpicture}[
        phaselabel/.style={font=\fignodefont, text=ggunavy, anchor=east},
        milestonelabel/.style={font=\figtitlefont, text=ggunavy, anchor=south, align=center},
        weeklabel/.style={font=\fignodefont, text=ggunavy}
    ]
        % --- Week axis ---
        \draw[thick, ggunavy] (0, 0) -- (13, 0);
        \foreach \w in {0,2,...,26} {
            \draw[ggunavy] ({\w*0.5}, -0.1) -- ({\w*0.5}, 0.1);
            \node[weeklabel, below] at ({\w*0.5}, -0.15) {\w};
        }
        \node[font=\figtitlefont, text=ggunavy, below] at (6.5, -0.55) {Week};

        % --- Phase bars (bottom to top) ---
        \fill[lightblue, rounded corners=2pt] ({1*0.5}, 0.4) rectangle ({4*0.5}, 0.7);
        \node[phaselabel] at (-0.1, 0.55) {Literature Review};

        \fill[lightorange, rounded corners=2pt] ({2*0.5}, 0.9) rectangle ({6*0.5}, 1.2);
        \node[phaselabel] at (-0.1, 1.05) {Dataset Acquisition};

        \fill[lightteal, rounded corners=2pt] ({5*0.5}, 1.4) rectangle ({8*0.5}, 1.7);
        \node[phaselabel] at (-0.1, 1.55) {Preprocessing Pipeline};

        \fill[lightpurple, rounded corners=2pt] ({7*0.5}, 1.9) rectangle ({13*0.5}, 2.2);
        \node[phaselabel] at (-0.1, 2.05) {Model Development};

        \fill[lightgreen, rounded corners=2pt] ({9*0.5}, 2.4) rectangle ({15*0.5}, 2.7);
        \node[phaselabel] at (-0.1, 2.55) {Training \& Optimization};

        \fill[lightyellow, rounded corners=2pt] ({14*0.5}, 2.9) rectangle ({21*0.5}, 3.2);
        \node[phaselabel] at (-0.1, 3.05) {Evaluation \& Testing};

        \fill[lightpink, rounded corners=2pt] ({19*0.5}, 3.4) rectangle ({23*0.5}, 3.7);
        \node[phaselabel] at (-0.1, 3.55) {Clinical Validation};

        \fill[lightcoral, rounded corners=2pt] ({22*0.5}, 3.9) rectangle ({26*0.5}, 4.2);
        \node[phaselabel] at (-0.1, 4.05) {Documentation};

        % --- Milestone markers (dashed vertical lines) ---
        \draw[dashed, thick, ggunavy!70] ({7*0.5}, -0.1) -- ({7*0.5}, 4.5);
        \node[milestonelabel] at ({7*0.5}, 4.6) {M1\\Data Ready};

        \draw[dashed, thick, ggunavy!70] ({14*0.5}, -0.1) -- ({14*0.5}, 4.5);
        \node[milestonelabel] at ({14*0.5}, 4.6) {M2\\Models Complete};

        \draw[dashed, thick, ggunavy!70] ({22*0.5}, -0.1) -- ({22*0.5}, 4.5);
        \node[milestonelabel] at ({22*0.5}, 4.6) {M3\\Evaluation Done};

        \draw[dashed, thick, ggunavy!70] ({26*0.5}, -0.1) -- ({26*0.5}, 4.5);
        \node[milestonelabel] at ({26*0.5}, 4.6) {M4\\Final};
    \end{tikzpicture}%
    }% end resizebox
    \caption{Research Timeline: 26-Week Project Plan with Milestones}
    \label{fig:research_timeline_appx}
    \vspace{0.5em}\noindent\textit{Note.} M = milestone. Overlapping phases indicate parallel work streams. M1 = data readiness; M2 = model architecture completion; M3 = evaluation and statistical testing; M4 = final submission.
\end{figure}


%% =====================================================================
\section{ASD-Focused Feature Extraction Comparison}
\label{sec:appx_feature_extraction}
%% =====================================================================

Table~\ref{tab:feature_comparison_appx} compares the feature extraction methods used across the five clinical domains, and Table~\ref{tab:signal_to_image_appx} documents the signal-to-image conversion methods that enable deep learning architectures to process time-series EEG data.

\needspace{8\baselineskip}
\iffalse % AUTO-SUPPRESS non-ASD
\needspace{8\baselineskip}
\noindent Table~\ref{tab:feature_comparison_appx} asd-focused feature extraction comparison.

{\tablestripes\tablefontsize
\begin{xltabular}{\textwidth}{@{} p{1.8cm} L L p{2cm} L L @{}}
\caption{ASD-Focused Feature Extraction Comparison}\label{tab:feature_comparison_appx} \\
\toprule
\thead{Feature Domain} & \thead{ASD} & \thead{PD} & \thead{Cancer} & \thead{Stress} & \thead{BCI} \\
\midrule
\endfirsthead
\multicolumn{6}{@{}l}{\textit{(\,Tab.~\ref{tab:feature_comparison_appx} continued from previous page\,)}} \\
\toprule
\thead{Feature Domain} & \thead{ASD} & \thead{PD} & \thead{Cancer} & \thead{Stress} & \thead{BCI} \\
\midrule
\endhead
\bottomrule
\endlastfoot
Time domain & Hjorth parameters, zero-crossing, statistical moments & Hjorth parameters, waveform shape analysis & N/A (pixel-level) & Statistical moments, Hjorth parameters & Statistical moments, zero-crossing rate \\
\addlinespace
Frequency & Band power ($\delta$, $\theta$, $\alpha$, $\beta$, $\gamma$), PSD, spectral entropy & Beta power (13--30 Hz), alpha power, $\theta$/$\beta$ ratio & N/A & Frontal $\alpha$ asymmetry, frontal $\theta$, high $\beta$ & Band power, PSD, spectral entropy \\
\addlinespace
Time-frequency & CWT scalogram & Multi-scale CWT (3 scales) & N/A & VMD + CWT & CWT spectrogram \\
\addlinespace
Spatial & Topographic maps & Channel attention maps & CNN feature maps (Conv1--5) & Channel-level features & Topographic maps \\
\addlinespace
Connectivity & PLV, coherence, Granger causality & Beta coherence, pathological synchrony & N/A & Limited (cross-channel) & Limited (cross-channel) \\
\addlinespace
Key biomarker & $\theta$/$\beta$ ratio elevation & Beta-band elevation and synchrony & Cell morphology (blast ratio) & Frontal $\alpha$ asymmetry & $\alpha$/$\theta$ ratio reduction \\
\end{xltabular}}
\vspace{0.5em}\noindent\textit{Note.} CWT = continuous wavelet transform; VMD = variational mode decomposition; PLV = phase-locking value; PSD = power spectral density; CNN = convolutional neural network; $\delta$ = delta (0.5--4 Hz); $\theta$ = theta (4--8 Hz); $\alpha$ = alpha (8--12 Hz); $\beta$ = beta (13--30 Hz); $\gamma$ = gamma ($>$30 Hz).
\fi % AUTO-SUPPRESS non-ASD

Table~\ref{tab:signal_to_image_appx} documents the signal-to-image conversion methods that enable deep learning architectures to process time-series EEG data across the five clinical domains.

\needspace{8\baselineskip}
\iffalse % AUTO-SUPPRESS non-ASD
\needspace{8\baselineskip}
{\tablestripes\tablefontsize
\begin{xltabular}{\textwidth}{@{} p{1.6cm} p{3.5cm} p{2.8cm} L @{}}
\caption{Signal-to-Image Conversion Methods Across Domains}\label{tab:signal_to_image_appx} \\
\toprule
\thead{Domain} & \thead{Conversion Method} & \thead{Output Dimensions} & \thead{Rationale} \\
\midrule
\endfirsthead
\multicolumn{4}{@{}l}{\textit{(\,Tab.~\ref{tab:signal_to_image_appx} continued from previous page\,)}} \\
\toprule
\thead{Domain} & \thead{Conversion Method} & \thead{Output Dimensions} & \thead{Rationale} \\
\midrule
\endhead
\bottomrule
\endlastfoot
ASD (EEG) & CWT scalogram (Morlet wavelet) & 224 $\times$ 224 $\times$ 3 (RGB) & Captures multi-scale oscillatory patterns; compatible with ImageNet-pretrained CNNs \\
\addlinespace
& Multi-scale CWT (3 frequency scales) & 3 $\times$ 64 $\times$ 64 & Separates $\alpha$, $\beta$, $\gamma$ band contributions for channel attention mechanism \\
\addlinespace
Cancer (Image) & Direct input (already image) & 224 $\times$ 224 $\times$ 3 (RGB) & Microscopy images require no signal conversion; stain normalisation applied instead \\
\addlinespace
Stress (EEG) & CWT + topographic map overlay & 128 $\times$ 128 $\times$ 3 (RGB) & Combines time-frequency and spatial activation patterns \\
\addlinespace
BCI (EEG) & Spectrogram (STFT) & 224 $\times$ 224 $\times$ 3 (RGB) & Standard short-time Fourier transform captures drowsiness-specific $\alpha$/$\theta$ transitions \\
\end{xltabular}}
\vspace{0.5em}\noindent\textit{Note.} CWT = continuous wavelet transform; STFT = short-time Fourier transform; CNN = convolutional neural network; RGB = red--green--blue colour channels.
\fi % AUTO-SUPPRESS non-ASD


%% =====================================================================
\section{Software Stack and Hardware Specifications}
\label{sec:appx_sw_hw}
%% =====================================================================

Table~\ref{tab:sw_stack_appx} documents the software stack and Table~\ref{tab:hw_reproducibility_appx} provides hardware specifications and reproducibility controls.
\needspace{8\baselineskip}
{\tablestripes\tablefontsize
\begin{xltabular}{\textwidth}{@{} p{2.5cm} p{4.5cm} L @{}}
\caption{Software Stack for Experimental Infrastructure}\label{tab:sw_stack_appx} \\
\toprule
\thead{Category} & \thead{Component} & \thead{Specification} \\
\midrule
\endfirsthead
\multicolumn{3}{@{}l}{\textit{(\,Tab.~\ref{tab:sw_stack_appx} continued from previous page\,)}} \\
\toprule
\thead{Category} & \thead{Component} & \thead{Specification} \\
\midrule
\endhead
\bottomrule
\endlastfoot
Core language & Python & 3.10+ \\
\addlinespace
Deep learning & PyTorch, TensorFlow, Keras & 2.0+, 2.x, integrated \\
\addlinespace
Signal processing & MNE-Python, SciPy, NumPy & Domain-specific EEG analysis \\
\addlinespace
Image processing & OpenCV, Scikit-learn & Radiomics and CV pipeline \\
\addlinespace
GenAI / RAG & LangChain, Hugging Face Transformers & RAG pipeline and embeddings \\
\addlinespace
GPU acceleration & CUDA, cuDNN & 11.8+, 8.6+ \\
\end{xltabular}}
\vspace{0.5em}\noindent\textit{Note.} CUDA = Compute Unified Device Architecture; cuDNN = CUDA Deep Neural Network library; RAG = retrieval-augmented generation; CV = computer vision.

Table~\ref{tab:hw_reproducibility_appx} summarises the hardware specifications and reproducibility controls employed across all experiments.
\needspace{8\baselineskip}
{\tablestripes\tablefontsize
\begin{xltabular}{\textwidth}{@{} p{2.5cm} p{4.5cm} L @{}}
\caption{Hardware Specifications and Reproducibility Controls}\label{tab:hw_reproducibility_appx} \\
\toprule
\thead{Category} & \thead{Component} & \thead{Specification} \\
\midrule
\endfirsthead
\multicolumn{3}{@{}l}{\textit{(\,Tab.~\ref{tab:hw_reproducibility_appx} continued from previous page\,)}} \\
\toprule
\thead{Category} & \thead{Component} & \thead{Specification} \\
\midrule
\endhead
\bottomrule
\endlastfoot
Primary GPU & NVIDIA RTX 4090 & 24 GB VRAM \\
\addlinespace
Alternative GPU & NVIDIA A100 & 40 GB VRAM \\
\addlinespace
CPU & AMD Ryzen 9 / Intel i9 & Multi-core processing \\
\addlinespace
Memory & System RAM & 64 GB DDR5 (minimum) \\
\addlinespace
Storage & NVMe SSD & 2 TB \\
\addlinespace
Training time & Per domain & 24--72 hours \\
\addlinespace
Environment & Docker containers & Pinned dependency versions \\
\addlinespace
Determinism & Fixed random seed (42) & PyTorch, TensorFlow, NumPy, Scikit-learn \\
\addlinespace
Version control & Git & Tagged experiment commits \\
\addlinespace
Experiment tracking & MLflow & Hyperparameter and metric logging \\
\end{xltabular}}
\vspace{0.5em}\noindent\textit{Note.} GPU = graphics processing unit; VRAM = video random access memory; DDR5 = Double Data Rate 5; NVMe = Non-Volatile Memory Express; SSD = solid-state drive; MLflow = open-source machine learning lifecycle platform.


%% =====================================================================
\section{GenAI Evaluation Frameworks Comparison}
\label{sec:appx_genai_eval}
%% =====================================================================

Table~\ref{tab:genai_eval_frameworks_appx} compares the three GenAI evaluation frameworks used in this study: DeepEval (primary evaluation harness), RAGAS (retrieval-specific metrics), and G-Eval (LLM-judged narrative quality).
\needspace{8\baselineskip}
{\tablestripes\tablefontsize
\begin{xltabular}{\textwidth}{@{} p{1.6cm} L L p{1.6cm} L @{}}
\caption{GenAI Evaluation Frameworks: DeepEval, RAGAS, and G-Eval}\label{tab:genai_eval_frameworks_appx} \\
\toprule
\thead{Framework} & \thead{Metrics} & \thead{Evaluation Method} & \thead{Open Source} & \thead{RGAIG Usage} \\
\midrule
\endfirsthead
\multicolumn{5}{@{}l}{\textit{(\,Tab.~\ref{tab:genai_eval_frameworks_appx} continued from previous page\,)}} \\
\toprule
\thead{Framework} & \thead{Metrics} & \thead{Evaluation Method} & \thead{Open Source} & \thead{RGAIG Usage} \\
\midrule
\endhead
\bottomrule
\endlastfoot
DeepEval & 14 metrics: faithfulness, hallucination, bias, toxicity, answer relevancy & Reference-free and reference-based; custom metric support & Yes (Python) & Primary framework for faithfulness and hallucination scoring \\
\addlinespace
RAGAS & 8 metrics: context precision, context recall, faithfulness, answer relevancy, answer correctness & Component-wise evaluation of retrieval and generation stages independently & Yes (Python) & Retrieval quality: context precision and recall measurement \\
\addlinespace
G-Eval & 4 criteria: coherence, consistency, fluency, relevance & LLM-based evaluation with chain-of-thought scoring (GPT-4 as judge) & Partial & Coherence and fluency scoring for clinical narrative quality \\
\end{xltabular}}
\vspace{0.5em}\noindent\textit{Note.} DeepEval serves as the primary evaluation harness; RAGAS provides retrieval-specific metrics; G-Eval provides LLM-judged narrative quality assessment. Using multiple frameworks provides triangulation of output quality.


%% =====================================================================
\section{RAG Output Quality Metrics}
\label{sec:appx_rag_metrics}
%% =====================================================================

Table~\ref{tab:dcaibf_output_metrics_appx} presents the RAG output quality metrics and acceptance thresholds applied during RGAIG validation testing across 500 generated explanations in five clinical domains.
\needspace{8\baselineskip}
{\tablestripes\tablefontsize
\begin{xltabular}{\textwidth}{@{} p{2cm} p{1.6cm} p{1cm} p{1.5cm} L @{}}
\caption{RGAIG RAG Output Quality Metrics and Acceptance Thresholds}\label{tab:dcaibf_output_metrics_appx} \\
\toprule
\thead{Metric} & \thead{Framework} & \thead{Score} & \thead{Threshold} & \thead{Interpretation} \\
\midrule
\endfirsthead
\multicolumn{5}{@{}l}{\textit{(\,Tab.~\ref{tab:dcaibf_output_metrics_appx} continued from previous page\,)}} \\
\toprule
\thead{Metric} & \thead{Framework} & \thead{Score} & \thead{Threshold} & \thead{Interpretation} \\
\midrule
\endhead
\bottomrule
\endlastfoot
Faithfulness & DeepEval & 0.89 & $\geq$0.85 & 89\% of generated claims supported by retrieved context \\
\addlinespace
Answer relevancy & RAGAS & 0.92 & $\geq$0.80 & Generated answers address the clinical query directly \\
\addlinespace
Context recall & RAGAS & 0.87 & $\geq$0.80 & 87\% of relevant passages retrieved from knowledge base \\
\addlinespace
Context precision & RAGAS & 0.91 & $\geq$0.85 & 91\% of retrieved passages are relevant to the query \\
\addlinespace
Hallucination rate & DeepEval & 2.1\% & $\leq$5\% & 2.1\% of claims have no supporting evidence \\
\addlinespace
Citation accuracy & Custom & 94.2\% & $\geq$90\% & 94.2\% of citations traceable to knowledge base \\
\addlinespace
Coherence & G-Eval & 4.6/5.0 & $\geq$4.0 & Expert-rated narrative quality and logical flow \\
\addlinespace
Toxicity & DeepEval & 0.0\% & $\leq$1\% & No toxic, biased, or harmful content detected \\
\end{xltabular}}
\vspace{0.5em}\noindent\textit{Note.} $n$=500 generated explanations across five clinical domains. Faithfulness = proportion of generated claims supported by retrieved context. Citation accuracy = proportion of cited references traceable to the knowledge base. Hallucination rate = percentage of generated claims with no supporting evidence.


%% =====================================================================
\section{End-to-End System Architecture: Eight Processing Flows}
\label{sec:appx_system_flows}
%% =====================================================================

Table~\ref{tab:system_flows_summary_appx} summarises the eight processing flows that comprise the RGAIG pipeline, documenting key stages, quality gates, and output artefacts for each flow.
\needspace{8\baselineskip}
{\tablestripes\tablefontsize
\begin{xltabular}{\textwidth}{@{} p{0.6cm} p{1.8cm} L p{2.8cm} L @{}}
\caption{End-to-End System Architecture: Eight Processing Flows}\label{tab:system_flows_summary_appx} \\
\toprule
\thead{\#} & \thead{Flow} & \thead{Key Stages} & \thead{Quality Gate} & \thead{Output Artefact} \\
\midrule
\endfirsthead
\multicolumn{5}{@{}l}{\textit{(\,Tab.~\ref{tab:system_flows_summary_appx} continued from previous page\,)}} \\
\toprule
\thead{\#} & \thead{Flow} & \thead{Key Stages} & \thead{Quality Gate} & \thead{Output Artefact} \\
\midrule
\endhead
\bottomrule
\endlastfoot
1 & Data Flow & Ingestion $\to$ Validation $\to$ Preprocessing $\to$ Feature extraction & SNR $>$10~dB; completeness $>$95\% & 142-feature vector per sample \\
\addlinespace
2 & Model Training Flow & Stratified split $\to$ 10-fold CV $\to$ baseline + DL training $\to$ Ensemble & All folds converge; no data leakage & Best model per domain \\
\addlinespace
3 & Hybrid Baseline+DL Flow & Dual-branch (3 baseline + 3 DL) $\to$ Majority-vote ensemble & Ensemble $>$ individual models & Final prediction with confidence \\
\addlinespace
4 & RAG Explainability Flow & Prediction $\to$ Query $\to$ FAISS retrieval $\to$ LLM generation $\to$ Hallucination check & Citation accuracy $\geq$90\%; coherence $\geq$4/5 & Evidence-grounded clinical narrative \\
\addlinespace
5 & GPU/vLLM Deployment & Model export $\to$ Quantisation $\to$ vLLM serving $\to$ API gateway & Latency $<$200~ms; accuracy drop $<$1\% & Production REST API endpoint \\
\addlinespace
6 & PII Protection Flow & PII detection $\to$ De-identification $\to$ Encryption $\to$ Access control $\to$ Audit & $k$-anonymity $\geq$5; zero PII in model input & HIPAA/GDPR-compliant data pipeline \\
\addlinespace
7 & Guardrail AI Flow & Input validation $\to$ Content filter $\to$ Hallucination detection $\to$ Confidence threshold & All 6 gates pass; else human escalation & Clinically safe, verified output \\
\addlinespace
8 & Statistical Flow & Hypothesis $\to$ Test selection $\to$ CV execution $\to$ Bootstrap CI $\to$ Effect size $\to$ Bonferroni & $p < .05$; CI excludes null; $d > 0.5$ & Accept/reject verdict per hypothesis \\
\end{xltabular}}
\vspace{0.5em}\noindent\textit{Note.} Each flow operates autonomously but shares data through structured JSON messages via the multi-agent orchestration layer. Quality gates enforce fail-fast semantics: a stage that fails its gate triggers retry or human escalation, never silent pass-through. vLLM = virtual large language model serving engine; FAISS = Facebook AI Similarity Search; PII = personally identifiable information.


%% =====================================================================
\section{Model Control Protocol: Governance Capabilities}
\label{sec:appx_mcp}
%% =====================================================================

Table~\ref{tab:mcp_capabilities_appx} documents the six governance capabilities implemented by the Model Control Protocol (MCP), which provides a standardised governance layer controlling how AI agents access tools, data sources, and models within the RGAIG pipeline.
\needspace{8\baselineskip}
{\tablestripes\tablefontsize
\begin{xltabular}{\textwidth}{@{} p{2cm} L L p{3.2cm} @{}}
\caption{Model Control Protocol: Six Governance Capabilities}\label{tab:mcp_capabilities_appx} \\
\toprule
\thead{Capability} & \thead{Function} & \thead{Implementation} & \thead{Enforcement} \\
\midrule
\endfirsthead
\multicolumn{4}{@{}l}{\textit{(\,Tab.~\ref{tab:mcp_capabilities_appx} continued from previous page\,)}} \\
\toprule
\thead{Capability} & \thead{Function} & \thead{Implementation} & \thead{Enforcement} \\
\midrule
\endhead
\bottomrule
\endlastfoot
Model registry & Central catalogue of all deployed models with version, hash, and metadata & MLflow model registry; SHA-256 integrity hash per model artefact & No model loads without registry entry \\
\addlinespace
Version control & Track model lineage: data version $\to$ feature version $\to$ model version $\to$ deployment & Git-based versioning; DVC for data and model artefacts & Deployment blocked if version chain incomplete \\
\addlinespace
Audit logs & Immutable record of all model predictions, inputs, outputs, and configuration & Append-only SQLite log; daily rotation; 90-day retention & Every inference call logged; no silent execution \\
\addlinespace
Access control & Role-based permissions: who can train, deploy, query, or rollback models & RBAC policy engine; three roles (admin, clinician, patient) & API key + role verification per request \\
\addlinespace
Deployment gates & Pre-deployment quality checks: accuracy $\geq$85\%, fairness variance $\leq$2\%, latency $<$200\,ms & Automated gate in CI/CD pipeline; all gates must pass & Failed gate blocks deployment; alerts to admin \\
\addlinespace
Rollback & Instant reversion to previous model version if post-deployment metrics degrade & Blue-green deployment; previous version kept hot & Auto-rollback if accuracy drops $>$3\% in 24h \\
\end{xltabular}}
\vspace{0.5em}\noindent\textit{Note.} MCP = Model Control Protocol; RBAC = role-based access control; DVC = Data Version Control; CI/CD = continuous integration/continuous deployment.


%% =====================================================================
\section{Input Guardrail Rules}
\label{sec:appx_input_guardrails}
%% =====================================================================

Table~\ref{tab:ch3_input_guardrails_appx} documents the five pre-processing safety checks applied before any query reaches the RAG pipeline.

\needspace{4\baselineskip}
\iffalse % AUTO-SUPPRESS non-ASD
\needspace{8\baselineskip}
{\tablestripes\tablefontsize
\begin{xltabular}{\textwidth}{@{} p{2cm} L L p{3.5cm} @{}}
\caption{Input Guardrail Rules: Five Pre-Processing Safety Checks}\label{tab:ch3_input_guardrails_appx} \\
\toprule
\thead{Rule} & \thead{Trigger Condition} & \thead{Action} & \thead{Example} \\
\midrule
\endfirsthead
\multicolumn{4}{@{}l}{\textit{(\,Tab.~\ref{tab:ch3_input_guardrails_appx} continued from previous page\,)}} \\
\toprule
\thead{Rule} & \thead{Trigger Condition} & \thead{Action} & \thead{Example} \\
\midrule
\endhead
\bottomrule
\endlastfoot
Block diagnosis intent & Query requests definitive diagnosis (``do I have cancer?'') & Block; redirect to clinician & ``Please consult your healthcare provider for diagnosis'' \\
\addlinespace
Block prescription & Query requests medication or dosage (``what drug should I take?'') & Block; display disclaimer & ``RGAIG does not provide prescription advice'' \\
\addlinespace
Redirect high-risk & Query mentions self-harm, suicide, or emergency symptoms & Redirect to crisis hotline + human clinician & Display crisis resources (988, emergency contacts) \\
\addlinespace
Refuse out-of-scope & Query outside the five clinical domains (ASD, PD, Cancer, Stress, BCI) & Refuse with explanation & ``This system is designed for neurodegenerative conditions'' \\
\addlinespace
Redact PII & Query contains patient name, DOB, SSN, or medical record number & Redact before processing; log detection event & ``John Smith'' $\to$ ``[PATIENT]'' \\
\end{xltabular}}
\vspace{0.5em}\noindent\textit{Note.} PII = personally identifiable information; DOB = date of birth; SSN = Social Security number. Rules are evaluated in priority order; the first matching rule terminates processing.
\fi % AUTO-SUPPRESS non-ASD


%% =====================================================================
\section{Output Guardrail Rules}
\label{sec:appx_output_guardrails}
%% =====================================================================

Table~\ref{tab:ch3_output_guardrails_appx} documents the five post-generation safety checks applied after LLM generation and before display to the user.
\needspace{8\baselineskip}
{\tablestripes\tablefontsize
\begin{xltabular}{\textwidth}{@{} p{2.2cm} L L p{2.8cm} @{}}
\caption{Output Guardrail Rules: Five Post-Generation Safety Checks}\label{tab:ch3_output_guardrails_appx} \\
\toprule
\thead{Rule} & \thead{Detection Method} & \thead{Action} & \thead{Threshold} \\
\midrule
\endfirsthead
\multicolumn{4}{@{}l}{\textit{(\,Tab.~\ref{tab:ch3_output_guardrails_appx} continued from previous page\,)}} \\
\toprule
\thead{Rule} & \thead{Detection Method} & \thead{Action} & \thead{Threshold} \\
\midrule
\endhead
\bottomrule
\endlastfoot
Downgrade medical authority & Tone classifier detects authoritative medical claims & Inject hedging language (``this may indicate,'' ``consider consulting'') & Authority score $>$0.7 \\
\addlinespace
Inject uncertainty & Confidence below threshold for classification & Add explicit uncertainty statement to output & Model confidence $<$85\% \\
\addlinespace
Block unsupported claims & Claim-to-evidence mapping finds ungrounded assertions & Remove claim; replace with ``insufficient evidence'' & Zero matching passages \\
\addlinespace
Hallucination gate & Faithfulness score below threshold & Replace output with ``I don't have enough information'' & Faithfulness $<$0.70 \\
\addlinespace
Human review escalation & Sensitive topic detected (end-of-life, paediatric, mental health) & Flag for clinician review before patient display & Topic classifier confidence $>$0.8 \\
\end{xltabular}}
\vspace{0.5em}\noindent\textit{Note.} Output rules are evaluated sequentially; multiple rules may apply to the same output. Human review escalation is the final safety net before content reaches the patient.


%% =====================================================================
\section{PII Protection Pipeline}
\label{sec:appx_pii_pipeline}
%% =====================================================================

Table~\ref{tab:pii_pipeline_appx} documents the six-stage PII protection pipeline implemented to ensure patient data protection throughout the RGAIG system.
\needspace{8\baselineskip}
{\tablestripes\tablefontsize
\begin{xltabular}{\textwidth}{@{} p{0.6cm} L L p{3.8cm} @{}}
\caption{PII Protection Pipeline: Six-Stage Implementation}\label{tab:pii_pipeline_appx} \\
\toprule
\thead{\#} & \thead{Process} & \thead{Technology} & \thead{Verification} \\
\midrule
\endfirsthead
\multicolumn{4}{@{}l}{\textit{(\,Tab.~\ref{tab:pii_pipeline_appx} continued from previous page\,)}} \\
\toprule
\thead{\#} & \thead{Process} & \thead{Technology} & \thead{Verification} \\
\midrule
\endhead
\bottomrule
\endlastfoot
1 & Detect: identify PII entities in input text & Regex patterns (SSN, DOB, email, phone) + spaCy NER & Recall $>$99\% on synthetic PII test set \\
\addlinespace
2 & Classify: categorise entities by sensitivity & PHI (HIPAA-defined), PII (general), sensitive (clinical notes) & Three-tier classification \\
\addlinespace
3 & Redact: replace entities with type tokens & ``John Smith'' $\to$ ``[PERSON]''; ``123-45-6789'' $\to$ ``[SSN]'' & Redacted text reviewed by audit layer \\
\addlinespace
4 & Encrypt: protect data at rest and in transit & AES-256-GCM at rest; TLS 1.3 in transit & FIPS 140-2 compliant; 90-day key rotation \\
\addlinespace
5 & Access control: restrict by role and purpose & RBAC with three roles; GDPR Art.~5(1)(b) purpose limitation & Access denied without valid role + purpose \\
\addlinespace
6 & Audit: immutable access and processing log & Append-only SQLite log; tamper-evident hash chain & HIPAA \S164.312(b) compliance \\
\end{xltabular}}
\vspace{0.5em}\noindent\textit{Note.} PHI = protected health information; PII = personally identifiable information; NER = named entity recognition; RBAC = role-based access control; GCM = Galois/Counter Mode; FIPS = Federal Information Processing Standards.


%% =====================================================================
\section{Security Framework Alignment}
\label{sec:appx_security}
%% =====================================================================

Table~\ref{tab:security_alignment_appx} maps RGAIG security controls to five regulatory frameworks, demonstrating compliance readiness.
\needspace{8\baselineskip}
{\tablestripes\tablefontsize
\begin{xltabular}{\textwidth}{@{} p{2.5cm} L L @{}}
\caption{Security Framework Alignment: Five Regulatory Standards}\label{tab:security_alignment_appx} \\
\toprule
\thead{Regulation} & \thead{Key Requirement} & \thead{RGAIG Control} \\
\midrule
\endfirsthead
\multicolumn{3}{@{}l}{\textit{(\,Tab.~\ref{tab:security_alignment_appx} continued from previous page\,)}} \\
\toprule
\thead{Regulation} & \thead{Key Requirement} & \thead{RGAIG Control} \\
\midrule
\endhead
\bottomrule
\endlastfoot
HIPAA \S164.312 & Technical safeguards: access control, audit, encryption, integrity & RBAC, AES-256, audit trail, hash verification \\
\addlinespace
GDPR Art.~25/32 & Data protection by design; encryption; pseudonymisation & Privacy-by-design architecture; $k$-anonymity $\geq$5 \\
\addlinespace
NIST AI RMF & Risk management throughout AI lifecycle; trustworthy AI & 15 governance layers (A1--A15); continuous monitoring \\
\addlinespace
ISO 27001 & Information security management system; risk assessment & Documented risk register; annual review cycle \\
\addlinespace
EU AI Act Art.~14 & Human oversight for high-risk AI systems & Human-in-the-loop escalation; confidence thresholds \\
\end{xltabular}}
\vspace{0.5em}\noindent\textit{Note.} HIPAA = Health Insurance Portability and Accountability Act; GDPR = General Data Protection Regulation; NIST AI RMF = National Institute of Standards and Technology AI Risk Management Framework; EU AI Act = European Union Artificial Intelligence Act.


%% =====================================================================
\section{Unified Quality Taxonomy}
\label{sec:appx_quality_taxonomy}
%% =====================================================================

Table~\ref{tab:unified_taxonomy_appx} presents the RGAIG unified quality taxonomy, comprising 525 modules across three tiers covering the full stack of biomedical AI quality assurance.
\needspace{8\baselineskip}
{\tablestripes\tablefontsize
\begin{xltabular}{\textwidth}{@{} p{2.8cm} p{2cm} C C X @{}}
\caption{RGAIG Unified Quality Taxonomy}\label{tab:unified_taxonomy_appx} \\
\toprule
\thead{Tier} & \thead{Section} & \thead{Phases/ Pillars} & \thead{Modules} & \thead{Validation Focus} \\
\midrule
\endfirsthead
\multicolumn{5}{@{}l}{\textit{(\,Tab.~\ref{tab:unified_taxonomy_appx} continued from previous page\,)}} \\
\toprule
\thead{Tier} & \thead{Section} & \thead{Phases/ Pillars} & \thead{Modules} & \thead{Validation Focus} \\
\midrule
\endhead
\bottomrule
\endlastfoot
Tier 1: DL Pipeline & Ch.~\ref{chap:methodology} & 8 & 111 & Data integrity, model selection, performance, subject reliability, sensitivity, statistics, benchmarking, monitoring \\
\addlinespace
Tier 2: GenAI/RAG Pipeline & Ch.~\ref{chap:methodology} & 13 & 220 & Knowledge quality, retrieval fidelity, generation grounding, agent behaviour, A2A coordination, MCP enforcement, explainability, robustness \\
\addlinespace
Tier 3: Enterprise Governance & Ch.~\ref{chap:methodology} & 10 & 194 & Energy efficiency, hallucination prevention, hypothesis rigour, threat modelling, SWOT, governance, compliance, responsible AI, explainability, security \\
\midrule
\textbf{Total} & & \textbf{31} & \textbf{525} & \textbf{Full-stack biomedical AI quality assurance} \\
\end{xltabular}}
\vspace{0.5em}\noindent\textit{Note.} Execution across 500 validation queries achieved 98.1\% mean pass rate for Tier~2. Tier~1 and Tier~3 results are reported in Chapter~\ref{chap:analysis}.


%% =====================================================================
\section{Digital Transformation Guidelines}
\label{sec:appx_digital_transformation}
%% =====================================================================

This section situates the RGAIG methodology within the broader digital transformation landscape of healthcare, connecting methodological design decisions to digital transformation use cases, measurable impact, and value creation.

Table~\ref{tab:dt_guidelines_appx} maps ten digital transformation guidelines from methodological design decisions to measurable clinical value.

\needspace{3\baselineskip}
\tablefontsize
\begin{xltabular}
{\textwidth}{@{} p{1.8cm} L p{3cm} L L @{}}
\caption{Digital Transformation Guidelines: From Methodology to Clinical Value} \label{tab:dt_guidelines_appx} \\
\toprule
\thead{Guideline} & \thead{Use Case} & \thead{Methodological Basis} & \thead{Impact} & \thead{Value Created} \\
\midrule
\endfirsthead
\endhead

\midrule
\multicolumn{5}{r}{\textit{Continued on next page}} \\
\endfoot
\bottomrule
\endlastfoot

DT-1: ASD-focused diagnostic pipeline & Replace fragmented, domain-specific diagnostic tools with a single ASD-focused AI framework & Design science research; ASD-focused preprocessing; ASD-focused feature engineering & Eliminates tool sprawl; reduces vendor lock-in; standardises diagnostic quality across departments & 15--25\% reduction in per-diagnosis infrastructure cost; single governance framework for all AI tools \\
\addlinespace

DT-2: Automated feature extraction & Replace manual EEG annotation and imaging feature measurement with automated DL pipelines & Eight-step EEG pipeline; imaging radiomics pipeline; 7 feature categories & Reduces diagnostic preparation time by 60--70\%; eliminates inter-rater variability in feature extraction & Clinician time freed for patient interaction; consistent feature quality regardless of operator experience \\
\addlinespace

DT-3: Explainable AI integration & Provide clinicians with interpretable, literature-grounded diagnostic explanations alongside predictions & SHAP feature importance + RAG narrative generation; expert panel validation & Transforms AI from black box to decision support; increases clinician trust and regulatory compliance & Higher adoption rates; reduced liability risk; FDA/EU AI Act readiness \\
\addlinespace

DT-4: Multi-agent workflow automation & Automate the end-to-end diagnostic workflow from data ingestion to report generation & Automated pipeline orchestration; domain-specific expert agents; coordinator agent & 60\% reduction in manual workflow steps; automated quality checks at each pipeline stage & Faster time-to-diagnosis; reduced human error in data handling; scalable to multiple departments \\
\addlinespace

DT-5: Governance-by-design & Embed AI governance, accountability, and audit mechanisms into the framework architecture rather than retrofitting & RACI matrix; escalation protocols; RTAI framework; data governance model & Proactive regulatory compliance; continuous monitoring and accountability; audit-ready from deployment & Avoided regulatory penalties; faster time-to-approval; institutional trust in AI systems \\
\addlinespace

DT-6: ASD-focused knowledge transfer & Enable diagnostic models trained on one biomedical domain to inform and improve performance in related domains & Multi-scale channel attention (ASD ensemble); shared feature space; transfer learning architecture & Reduces per-ASD training data requirements by 20--30\%; accelerates deployment to new clinical areas & Faster expansion to new specialities; reduced data collection burden for rare conditions \\
\addlinespace

DT-7: Privacy-preserving data pipelines & Process sensitive patient data without exposing personally identifiable information & De-identification protocols; HIPAA/GDPR compliance; public dataset benchmarking & Full regulatory compliance; patient trust preservation; institutional risk mitigation & Zero privacy violations; compatible with federated learning (future deployment) \\
\addlinespace

DT-8: Evidence-based decision support & Ground every AI recommendation in quantifiable evidence with statistical confidence bounds & Bootstrap resampling (1,000$\times$); 95\% CI; effect sizes (Cohen's $d$); predefined thresholds & Clinicians receive confidence intervals, not point estimates; uncertainty is quantified and communicated & Better-informed clinical decisions; reduced over-reliance on single AI predictions \\
\addlinespace

DT-9: Scalable deployment architecture & Design the framework for institutional scale-up from pilot to enterprise-wide deployment & Modular pipeline design; containerisable components; standardised APIs; batch and real-time modes & Single framework scales from 1 department to hospital-wide deployment without architectural redesign & Reduced total cost of ownership; linear scaling of operational benefits \\
\addlinespace

DT-10: Continuous improvement loop & Enable the AI system to improve over time through structured feedback, retraining, and performance monitoring & Validation framework; predefined reliability thresholds; drift detection design & AI accuracy maintains or improves post-deployment; degradation is detected and corrected automatically & Long-term ROI; institutional learning; competitive advantage through AI maturity \\

\end{xltabular}

Table~\ref{tab:dt_impact_quantification_appx} quantifies the expected digital transformation impact across five measurable dimensions.
\needspace{3\baselineskip}
{\tablestripes\tablefontsize
\begin{xltabular}{\textwidth}{@{} p{2.2cm} p{2cm} L p{2cm} p{2.5cm} @{}}
\caption{Digital Transformation Impact Quantification}\label{tab:dt_impact_quantification_appx} \\
\toprule
\thead{Impact Dimension} & \thead{Baseline (No AI)} & \thead{With RGAIG} & \thead{Improvement} & \thead{Evidence Source} \\
\midrule
\endfirsthead
\multicolumn{5}{@{}l}{\textit{(\,Tab.~\ref{tab:dt_impact_quantification_appx} continued from previous page\,)}} \\
\toprule
\thead{Impact Dimension} & \thead{Baseline (No AI)} & \thead{With RGAIG} & \thead{Improvement} & \thead{Evidence Source} \\
\midrule
\endhead
\bottomrule
\endlastfoot
Diagnostic accuracy & 78--82\% (manual clinician) & 97.67\% $\pm$ 2.6\% (ASD ensemble) & +8--12 pp & Ch.~\ref{chap:analysis}; AUC = 0.956 \\
\addlinespace
Diagnostic time per case & 25--35 minutes & 10--15 minutes (automated pipeline) & $-$55--65\% & 60\% manual effort reduction (H4) \\
\addlinespace
Workflow automation & 0\% automated (fully manual) & 60--70\% automated (automated agents) & +60--70 pp & H4 validation: coordinator + domain agents handle preprocessing, classification, explanation \\
\addlinespace
Explainability compliance & Ad hoc or absent & SHAP + RAG narrative per diagnosis & 0\% $\rightarrow$ 100\% coverage & H5 validation: expert panel mean $\geq$ 4.0/5.0; Krippendorff's $\alpha = 0.84$ \\
\addlinespace
Governance readiness & Fragmented or absent & RACI + escalation + KPI monitoring & 0\% $\rightarrow$ full coverage & RTAI framework: 46 modules, 5 pillars, 77\% maturity \\
\end{xltabular}}
\vspace{0.5em}\noindent\textit{Note.} pp = percentage points; ASD ensemble = multi-scale channel attention; AUC = area under the receiver operating characteristic curve. Baseline figures are drawn from published healthcare diagnostic literature.

Table~\ref{tab:dt_value_framework_appx} maps the digital transformation value creation to four stakeholder perspectives, following a balanced scorecard structure.
\needspace{3\baselineskip}
{\tablestripes\tablefontsize
\begin{xltabular}{\textwidth}{@{} p{2cm} L L p{2.5cm} @{}}
\caption{Digital Transformation Value Framework: Four Stakeholder Perspectives}\label{tab:dt_value_framework_appx} \\
\toprule
\thead{Perspective} & \thead{Value Created by RGAIG} & \thead{Key Metrics} & \thead{Supporting Guidelines} \\
\midrule
\endfirsthead
\multicolumn{4}{@{}l}{\textit{(\,Tab.~\ref{tab:dt_value_framework_appx} continued from previous page\,)}} \\
\toprule
\thead{Perspective} & \thead{Value Created by RGAIG} & \thead{Key Metrics} & \thead{Supporting Guidelines} \\
\midrule
\endhead
\bottomrule
\endlastfoot
Patient & Faster, more accurate diagnoses; reduced misdiagnosis risk; transparent AI explanations for informed consent & Diagnostic accuracy (+8--12 pp); false negative reduction ($-$50\%); time-to-diagnosis ($-$55\%) & DT-1, DT-3, DT-8 \\
\addlinespace
Clinician & Reduced cognitive workload; decision support with confidence intervals; automated routine tasks; professional development through AI collaboration & Cognitive workload ($-$0.8 points); cases per shift (+40--60\%); satisfaction (+0.8--1.2 points) & DT-2, DT-3, DT-4, DT-8 \\
\addlinespace
Administrator & Lower per-diagnosis cost; regulatory compliance; reduced liability; scalable infrastructure; competitive positioning & Cost per diagnosis ($-$15--25\%); governance coverage (0 $\rightarrow$ 100\%); vendor consolidation & DT-1, DT-5, DT-7, DT-9 \\
\addlinespace
IT / Informatics & Modular, maintainable architecture; standardised pipelines; automated monitoring; continuous improvement capability & Deployment time (weeks $\rightarrow$ days for new domains); monitoring coverage (100\%); drift detection & DT-6, DT-9, DT-10 \\
\end{xltabular}}
\vspace{0.5em}\noindent\textit{Note.} pp = percentage points. Metrics are drawn from the computational validation (Chapter~\ref{chap:analysis}), ROI analysis (Chapter~\ref{chap:discussion}), and Phase~2 pre/post performance impact assessment (planned). The balanced scorecard structure ensures digital transformation value is assessed from multiple stakeholder perspectives.

These digital transformation guidelines argue that the methodological design of the RGAIG is not merely an academic exercise but a deliberate engineering strategy for clinical deployment. Each guideline (DT-1 through DT-10) traces a direct line from a methodological decision to a measurable operational impact and a stakeholder-specific value proposition.


%% =====================================================================
\section{Comprehensive Assessment Framework Summary}
\label{sec:appx_assessment}
%% =====================================================================

The comprehensive assessment framework for Chapter~3---covering reliability and validity, methodological comparability with prior work, contribution claims (C1--C6), TAM methodology design, trust and algorithm aversion mitigation, qualitative--quantitative analysis balance, governance and accountability design, and contributions summary---confirms that the research methodology meets standards for internal validity (stratified 10-fold CV), reliability (bootstrap 1,000 resamples), statistical conclusion validity (power analysis, $n$=26, Bonferroni correction), and content validity (expert panel protocol, $n$=12).
